% source: 论文起草.md L9-L78
\section{问题重述}\label{sec:restate}

\subsection{问题背景}\label{sec:background}

全球 AI 算力需求正大幅增长，部分规划中的超大规模园区，电力需求已达到城市级量级。然而，电网容量、生态基流、土地刚性约束，使超大规模集中式园区受局部资源的上限限制；与此同时，AI 基础设施高度集中于少数热点地区，广大资源欠发达地区仍处于算力版图之外。

为解决这一问题，本文通过数学建模回答一个核心政策问题：\textbf{集中式 AI 数据中心的规模上限由什么因素决定？弹性分布式网络能否在环境代价、资源压力与普惠性之间找到更优解？}

本题的核心，是对“集中式扩张”与“弹性分布式协同”两条路线进行量化边界测算与优劣比较：

\subsection{问题一：集中式 AIDC 的最大可持续规模}\label{sec:problem1}

只要一项资源触及当地极限，园区就无法继续扩容，呈现典型的木桶效应。

因此，对给定地理场址与建设运行配置 $x$（供电、冷却与用地组合，且受电源建设工期制约），分别测算其在能源、水资源和土地三类约束下可长期维持的额定 IT 功率上限 $P_E(x)$、$P_W(x)$ 和 $P_L(x)$（单位：MW IT）。在可行配置集合 $\mathcal X$ 中，该场址的最大可持续规模为：

\begin{equation}\label{eq:p-capacity}
P(x)=\min\{P_E(x),P_W(x),P_L(x)\}
\end{equation}
\begin{equation}\label{eq:p-max}
P_{\max}(y)=\max_{x\in\mathcal X_y}P(x).
\end{equation}

在此基础上，将模型测算的上限量级与 950 MW、2,263 MW、5 GW 等现实项目对标以解释差异成因，并评估煤电、燃气、风电与光伏因水耗和占地不同而对上限带来的改变。

\subsection{问题二：弹性分布式算力网络的替代效益与边界}\label{sec:problem2}

在总算力需求相同的条件下，需要同时考察环境足迹总量与资源消耗的空间分布。分散部署可能降低单个地区的资源压力；与此同时，节点间通信和冷却耗能差异也可能带来额外消耗，增加总用电、耗水或占地。

为比较两种模式，在给定总算力需求 $O_{\mathrm{ut}}$ 下，我们引入环境足迹指数 $F$。以集中式方案的碳、水、地足迹 $C_0,W_0,A_0$ 为基准，对数加权衡量全系统总消耗的相对增减。

\begin{equation}\label{eq:p-f}
F=\ln (w_C\frac{C}{C_{\mathrm{ref}}}+w_W\frac{W}{W_{\mathrm{ref}}}+w_A\frac{A}{A_{\mathrm{ref}}})
\end{equation}

以及局地资源压力指数 $S$。用非线性惩罚函数 $\phi(R)=\frac{R}{1-R}$ 刻画各节点水、电、地占用情况。分布式取最紧张节点，且任意 $R\geq1$ 则判为不可行。构造综合代价得分 $\Omega$：

\begin{equation}\label{eq:p-s}
S=\sum_{k\in\{W,E,A\}}m_k\phi(R_k).
\end{equation}

\begin{equation}\label{eq:p-omega}
\Omega=F+\lambda S.
\end{equation}

其中 $\sum w=\sum m=1$，$\lambda$ 权衡局地减压与全局足迹的相对重要性。若 $\Omega_{\mathrm{dis}}<\Omega_{\mathrm{cen}}$，则分布式方案综合占优；由此可界定分布式网络成立的条件，并进一步讨论其提升边缘地区算力可及性，即普惠性对应的环境与网络效率代价。
